\section{Caracterización del proceso CIP}
{Un ciclo de limpieza CIP consiste típicamente en la circulación programada de distintos fluidos (agua, solución detergente y agua de enjuague) a través del equipo que se desea higienizar, sin necesidad de desmontarlo \cite{bowser2023cip} \cite{johansson2025cip}. Cada fase del ciclo cumple una función específica: el preenjuague retira los residuos gruesos y reduce la carga orgánica antes del contacto con el detergente; el lavado emplea una solución química, generalmente a temperatura elevada, para disolver y remover residuos adheridos a las superficies; y el enjuague final retira los remanentes de detergente y devuelve el equipo a una condición apta para su reutilización \cite{johansson2025cip}. Para este prototipo, el equipo optó por un esquema de un solo paso: en cada fase, el fluido se prepara en la cabina de la unidad portátil, una única bomba de impulsión lo envía hacia el equipo de prueba y, tras recorrer sus superficies internas, se conduce directamente al drenaje. El sistema cuenta con dos bombas de funciones diferenciadas: la bomba de impulsión, que transporta el fluido desde la cabina hasta el equipo de prueba, y la bomba dosificadora, que agrega el detergente a la cabina únicamente durante la fase de lavado. Esta decisión simplifica la arquitectura hidráulica del PMV y evita la acumulación de sólidos o residuos de detergente dentro del circuito de preparación, a costa de un mayor consumo de agua y detergente por ciclo, aspecto que el equipo documenta como una limitación de diseño de este primer corte.}
\subsection{Fases del ciclo}
La Tabla 2 resume la caracterización preliminar de las tres fases del ciclo CIP. Para cada fase se indican el fluido y su volumen, la temperatura, el tiempo de operación, el estado de las bombas y válvulas, la condición de transición y el destino del fluido. Los valores de temperatura y tiempo son preliminares y se validarán experimentalmente en el segundo corte.
\begin{table}[H]\centering\scriptsize\caption{Caracterización preliminar de las fases CIP.}
\begin{tabularx}{\textwidth}{L{1.9cm}L{2.3cm}L{1.4cm}L{2.0cm}YL{1.4cm}L{1.6cm}}\toprule
\textbf{Fase}&\textbf{Fluido}&\textbf{Temp.}&\textbf{Tiempo}&\textbf{Bomba/válvulas}&\textbf{Transición}&\textbf{Destino}\\\midrule
Preenjuague&Agua potable, 15--20 L, máx. 20 L (paso único)&50\,$^{\circ}$C (preliminar)&2--3 min (preliminar; se define según el tiempo de paso del líquido total por las tuberías)&Bomba de impulsión encendida; bomba dosificadora apagada; válvula de salida a drenaje abierta&Sin cambios&Drenaje\\
Lavado&Agua potable, 15--20 L, máx. 20 L, con detergente diluido al 1,5\,\% v/v, aprox. 225--300 mL (paso único)&40\,$^{\circ}$C (preliminar)&5--10 min (preliminar; se define según el tiempo de paso del líquido total por las tuberías)&Bomba de impulsión encendida; bomba dosificadora activa solo durante la dosificación; válvulas dirigen el fluido en un solo paso hacia el equipo de prueba y de ahí al drenaje&Sin cambios&Drenaje\\
Enjuague final&Agua potable, 15--20 L, máx. 20 L (paso único)&Ambiente&2--3 min (preliminar; se define según el tiempo de paso del líquido total por las tuberías)&Bomba de impulsión encendida; bomba dosificadora apagada; válvula de salida a drenaje abierta&Sin cambios&Drenaje\\\bottomrule
\end{tabularx}\end{table}

{Nota: el sistema opera en esquema de un solo paso: el fluido recorre el equipo de prueba una sola vez y se conduce en su totalidad al drenaje. Los valores de temperatura, tiempo y volumen se marcan como preliminares y quedan sujetos a validación experimental y a la ficha técnica del detergente seleccionado durante el segundo corte.}

\subsection{Clasificación de la información utilizada}
Para distinguir lo que está sustentado de lo que aún es una hipótesis, la Tabla 3 clasifica la información de diseño utilizada en cuatro categorías: información confirmada, requisitos definidos por el proyecto, supuestos adoptados para el diseño y aspectos por validar experimentalmente. Para cada elemento se indica la sección, el requisito o la tabla que lo respalda.
La Tabla~\ref{tab:clasificacion-info-diseno} clasifica la información de diseño utilizada en cuatro categorías: información confirmada, requisitos definidos por el proyecto, supuestos adoptados para el diseño y aspectos por validar experimentalmente. Para cada elemento se indica la sección, el requisito o la tabla que lo respalda.

\newlength{\anchoflex}
\setlength{\anchoflex}{\dimexpr\textwidth-2.6cm-2.7cm-3.5cm-8\tabcolsep\relax}

\begin{footnotesize}
\begin{longtable}{L{2.6cm}L{\anchoflex}L{2.7cm}L{3.5cm}}
\caption{Clasificación de la información de diseño.}\label{tab:clasificacion-info-diseno} \\
\toprule
\textbf{Elemento}&\textbf{Descripción o valor}&\textbf{Fuente o justificación}&\textbf{Clasificación}\\
\midrule
\endfirsthead

\multicolumn{4}{l}{\small \tablename\ \thetable{} -- continuación}\\
\toprule
\textbf{Elemento}&\textbf{Descripción o valor}&\textbf{Fuente o justificación}&\textbf{Clasificación}\\
\midrule
\endhead

\midrule
\multicolumn{4}{r}{\small Continúa en la siguiente página}\\
\endfoot

\bottomrule
\endlastfoot

Temperatura de preenjuague & 50\,°C & Tabla 2, definida por el equipo con base en la dinámica térmica del proceso & Supuesto (por validar en Corte 2)\\
Temperatura de lavado & 40\,°C & Tabla 2 y REQ13 (rango de control 40--50\,°C; error en estado estacionario $\leq 2$\,°C, equivalente al 5\,\% de la consigna de 40\,°C; ver Nota 1) & Requisito\\
Temperatura de enjuague final & Ambiente (sin calentamiento) & Tabla 2, ciclo III del diseño conceptual (Sección 4.1) & Confirmado\\
Umbral de seguridad de temperatura & 70\,°C (corte automático del calentador) & REQ14, límite térmico del termostato NC (Sección 5.3) & Requisito\\
Sensor de temperatura & RTD Pt100 de 3 hilos, clase B, con transmisor 4--20\,mA (0--100\,°C) conectado a entrada analógica del PLC & REQ12, matriz Pugh de sensores (Sección 3.3), presupuesto preliminar (Tabla 35) y Sección 5.1 (escalamiento NORM\_X/SCALE\_X en TIA Portal) & Confirmado\\
Potencia del calentador & Una (1) resistencia sumergible de 1500\,W, única fuente de calentamiento del sistema & Sección 5.2, variable manipulada (potencia nominal de 1500\,W) & Confirmado\\
Masa activa de fluido en el modelo térmico & 15\,kg & Sección 5.5, parámetro del modelo de primer orden ($\tau \approx 261{,}6$\,min) & Supuesto\\
Coeficiente global de pérdidas (UA) & $\approx 4$\,W/K & Sección 5.3, hipótesis más débil del modelo, pendiente de ensayo de enfriamiento libre & Supuesto (por validar)\\
Bomba de impulsión del fluido & Motobomba centrífuga de 0.5\,HP, autocebante; única bomba de impulsión del sistema & Presupuesto preliminar (Tabla 35), Sección 4.1 & Confirmado\\
Bomba dosificadora de detergente & Peristáltica de 24\,V DC con tubo Norprene, accionada por relé del PLC (sin PWM); única bomba de dosificación & REQ09, presupuesto preliminar (Tabla 35) & Confirmado\\
Controlador lógico programable & PLC Siemens S7-1200 CPU 1214C AC/DC/Rly & REQ13, prestado por la universidad (Tabla 35) & Confirmado\\
Protocolo de comunicaciones PLC--HMI & PROFINET sobre Ethernet industrial & Sección 6.3, comparación preliminar de protocolos (Tabla 27) & Requisito (decisión de diseño)\\
Protocolo de integración PLC--ESP32 (IoT) & Modbus TCP, con MQTT hacia el servidor (AWS) & Sección 6.3 & Requisito (decisión de diseño)\\
Peso del sistema en condición de vacío & $\leq 23$\,kg (referencia NIOSH, \textit{National Institute for Occupational Safety and Health}) & REQ03, requisito de portabilidad & Requisito\\
Capacidad de los tanques & Agua/mezcla $\leq 20$\,L, jabón $\leq 5$\,L, agua residual $\leq 20$\,L & REQ06; tanque de proceso de 20\,L en presupuesto (Tabla 35) & Requisito\\
Esquema hidráulico del ciclo & Un solo paso: todo el fluido se conduce al drenaje tras recorrer el equipo de prueba & Sección 2, nota de la Tabla 2 & Confirmado (decisión de diseño del equipo)\\

\end{longtable}
\end{footnotesize}

\noindent{\footnotesize \textbf{Nota 1:} la tolerancia de $\pm 2$\,°C es un requisito definido por el equipo, no un valor normativo de la industria CIP. Se deriva del criterio de referencia del enunciado del Proyecto Integrador (error en estado estacionario $\leq 5$\,\%), que para una consigna de 40\,°C equivale a 2\,°C. Es compatible con la instrumentación: un RTD Pt100 clase B según IEC 60751 presenta una tolerancia de $\pm(0{,}30 + 0{,}005\,|t|)$\,°C, es decir, aproximadamente $\pm 0{,}5$\,°C a 40\,°C~\cite{iec60751}, por lo que el sensor resuelve la banda con margen. Su validez se verificará experimentalmente en el segundo corte.}

\subsection{Secuencia preliminar}
La Figura 1 presenta la secuencia general del ciclo de mejora continua aplicado al proceso CIP, que organiza el desarrollo del proyecto en seis etapas: identificación, análisis, diseño, implementación, validación, y seguimiento y mejora. Las tres primeras corresponden a este primer corte; la implementación y la validación se abordarán en los cortes siguientes, y el seguimiento se realizará de manera continua.
\begin{figure}[H]
\centering
\begin{tikzpicture}[
  fase/.style={
    rectangle,
    rounded corners=3pt,
    draw=AzulUSTA,
    line width=0.7pt,
    fill=AzulUSTA!8,
    text width=0.72\textwidth,
    minimum height=1.25cm,
    align=center,
    inner sep=7pt
  },
  flecha/.style={-{Stealth[length=3mm,width=2mm]}, line width=0.8pt, draw=AzulUSTA}
]
  \node[fase] (identificacion) {\textbf{1. Identificación}\\
  \node[fase, below=8mm of identificacion] (analisis) {\textbf{2. Análisis}\\
  \node[fase, below=8mm of analisis] (diseno) {\textbf{3. Diseño}\\
  \node[fase, below=8mm of diseno] (implementacion) {\textbf{4. Implementación}\\
  \node[fase, below=8mm of implementacion] (validacion) {\textbf{5. Validación}\\
  \node[fase, below=8mm of validacion] (seguimiento) {\textbf{6. Seguimiento y mejora}\\

  \draw[flecha] (identificacion) -- (analisis);
  \draw[flecha] (analisis) -- (diseno);
  \draw[flecha] (diseno) -- (implementacion);
  \draw[flecha] (implementacion) -- (validacion);
  \draw[flecha] (validacion) -- (seguimiento);
  \draw[flecha] (seguimiento.east) to[out=0,in=0,looseness=1.2]
    node[right, align=left, font=\footnotesize\itshape, text=AzulUSTA]
    {Retroalimentación para la mejora continua} (identificacion.east);
\end{tikzpicture}
\caption{Secuencia general del ciclo de mejora continua aplicado al proceso CIP.}
\label{fig:secuencia-cip}
\end{figure}

{Respecto a las condiciones de operación del ciclo, se prevé que la transición entre fases se realice por tiempo programado en el PLC. Adicionalmente, podrá establecerse un umbral de nivel mínimo en el tanque principal que apague la bomba de impulsión para evitar su funcionamiento en vacío y habilite la fase siguiente. Ante una condición de sobretemperatura, bajo nivel o accionamiento de una parada de emergencia, el sistema pasará a un estado de parada segura que detenga el calentamiento y ambas bombas, mantenga cerradas las válvulas de suministro y permita el drenaje controlado del fluido en proceso. La pérdida de comunicación con la capa de supervisión no detendrá el control local: el PLC mantendrá el control térmico y la secuencia del ciclo, y la supervisión se restablecerá al recuperarse el enlace. La finalización normal del ciclo se dará al completarse el enjuague final a través de un temporizador.}
